% Chapter 6, section 4: Memory state
% Write in plain LaTeX. Code of a guide step: \guide{<step>}; commands: \begin{shell}...\end{shell};
% a screenshot: \screenshot{file}{caption}; a gap in the draft: \placeholder{...}.
% See book/latex/queen.sty for the full list.

\section{Memory state}

Each node keeps its whole state in memory: queues, partitions, consumer positions, KV, timers,
request ids, and the dedup records of recent pushes. It is loaded once, at boot.

Messages are not in memory. They stay in the queue log files and are read from there. So memory
grows with the number of partitions and keys, and with the pushes of the dedup window, not with how
many messages a queue keeps.

\queenfig{memory-disk}{What lives in memory, and what lives on disk.}

Message traces are kept apart, in their own file, and are read from disk.
